\chapter{XXX对XXXX的影响}\label{chap:ch4}

本段为章前导语，主要用于对本章研究内容进行简要引入与概括说明。通常可包括本章研究目的、主要研究内容、拟解决的问题以及章节结构安排等，以增强章节之间的衔接性和整体逻辑性。

\section{材料与方法}\label{sec:ch4-methods}

\subsection{材料与试剂}\label{subsec:ch4-materials}

本节用于填写试验原料来源、主要化学试剂名称、纯度等级、生产厂家及批次信息等。

\subsection{仪器与设备}\label{subsec:ch4-instruments}

本节用于填写主要实验仪器、型号、生产厂家及关键测试设备信息。

\subsection{XXXX的制备}\label{subsec:ch4-preparation}

本节用于描述样品制备、处理条件设置及工艺流程。

\subsection{XXXX的测定}\label{subsec:ch4-determination}

本节用于描述指标测定方法、计算公式和操作要点。

\subsection{统计与分析}\label{subsec:ch4-statistics}

统计与分析方法同 \ref{subsec:ch2-statistics} 节，如有特殊情况另作说明。

\section{结果与讨论}\label{sec:ch4-results}

\subsection{XXXXXXXXX}\label{subsec:ch4-result-a}

本部分用于展示研究结果，并结合相关理论与文献对结果进行分析讨论，阐明其变化规律、影响因素及可能作用机制。涉及上下标、公式、变量符号及统计学表达时，建议采用规范的 \LaTeX{} 数学格式书写，如$\mathrm{O}_2^{\cdot-}$、$k_{\mathrm{cat}}$、$R^2$、$p<0.05$、$\Delta E$、$\mu\mathrm{m}$、$10^{-3}$、$a_i^{(n)}$、$\mathrm{PPO}_{\mathrm{activity}}$等。

\section{小结}\label{sec:ch4-summary}

本节用于归纳本章主要结果，概括关键发现并为下一章研究奠定基础。